% PREAMBLE FOR FRACTAL RESONANCE MATHEMATICS TEXTBOOK
% Author: Pablo Cohen
% ORCID: 0009-0002-0734-5565

% ============================================
% PACKAGES
% ============================================

% Page layout and formatting
\usepackage[letterpaper,margin=1in,includefoot]{geometry}
\usepackage{fancyhdr}
\usepackage{setspace}

% Mathematics
\usepackage{amsmath,amssymb,amsthm}
\usepackage{mathtools}
\usepackage{physics}
\usepackage{tensor}

% Graphics and color
\usepackage{graphicx}
\usepackage[dvipsnames,table]{xcolor}
\usepackage{tikz}
\usetikzlibrary{arrows,shapes,positioning,patterns,calc,backgrounds}
\usepackage{pgfplots}
\pgfplotsset{compat=1.18}
\usepgfplotslibrary{fillbetween}

% Better float placement
\usepackage{float}
\usepackage{placeins}  % Provides \FloatBarrier
\renewcommand{\topfraction}{0.9}
\renewcommand{\bottomfraction}{0.8}
\setcounter{topnumber}{2}
\setcounter{bottomnumber}{2}
\setcounter{totalnumber}{4}
\renewcommand{\textfraction}{0.07}
\renewcommand{\floatpagefraction}{0.7}

% Principia Fractalis color palette
\definecolor{pfblack}{HTML}{000000}
\definecolor{pfgold}{HTML}{F0C419}
\definecolor{pfblue}{HTML}{4AA3FF}
\definecolor{pfwhite}{HTML}{FFFFFF}

% Tables and lists
\usepackage{booktabs}
\usepackage{multirow}
\usepackage{enumitem}
% Register single-token enumerate label shortcuts used across the
% manuscript: \begin{enumerate}[(i)], \begin{enumerate}[(a)], etc.
% \SetEnumitemKey is the documented enumitem API for declaring
% these aliases (cf.\ enumitem manual, §"Cloning the basic lists").
\SetEnumitemKey{(i)}{label=(\roman*)}
\SetEnumitemKey{(a)}{label=(\alph*)}
\SetEnumitemKey{(A)}{label=(\Alph*)}
\SetEnumitemKey{(I)}{label=(\Roman*)}
\SetEnumitemKey{(1)}{label=(\arabic*)}
\usepackage{longtable}  % multi-page tables: \endhead, \endfoot, etc.
\usepackage{array}      % extended column specifiers

% Code listings
\usepackage{listings}
\usepackage{algorithm}
\usepackage{algorithmic}

% References and citations
\usepackage{url}
\usepackage{hyperref}
\usepackage{cleveref}
\usepackage{cite}

% Boxes and frames
\usepackage{tcolorbox}
\tcbuselibrary{theorems,skins,breakable}

% Index
\usepackage{makeidx}
\makeindex

% Fonts and Unicode support
\usepackage[T1]{fontenc}
\usepackage[utf8]{inputenc}
\usepackage{lmodern}

% Define emoji replacements for pdflatex
\DeclareUnicodeCharacter{1F7E2}{[L1]} % Green circle
\DeclareUnicodeCharacter{1F7E1}{[L2]} % Yellow circle
\DeclareUnicodeCharacter{1F534}{[L3]} % Red circle
\DeclareUnicodeCharacter{1F4BB}{[CODE]} % Computer
\DeclareUnicodeCharacter{1F4CA}{[GRAPH]} % Bar chart
\DeclareUnicodeCharacter{1F3AF}{[TARGET]} % Target
\DeclareUnicodeCharacter{26A0}{[!]} % Warning sign
\DeclareUnicodeCharacter{039B}{$\Lambda$} % Lambda
\DeclareUnicodeCharacter{03A9}{$\Omega$} % Omega
\DeclareUnicodeCharacter{2713}{\checkmark} % Heavy check mark (used in version_history)
\DeclareUnicodeCharacter{FE0F}{} % Variation selector-16 (emoji presentation, render nothing)
\DeclareUnicodeCharacter{1F4D6}{[BOOK]} % Open book emoji
\DeclareUnicodeCharacter{1F52C}{[LAB]} % Microscope emoji
\DeclareUnicodeCharacter{1F50D}{[SEARCH]} % Magnifying glass
\DeclareUnicodeCharacter{1F4C8}{[CHART]} % Chart with upward trend
\DeclareUnicodeCharacter{1F4C9}{[CHART]} % Chart with downward trend
\DeclareUnicodeCharacter{1F4DD}{[NOTE]} % Memo / pencil
\DeclareUnicodeCharacter{1F511}{[KEY]} % Key
\DeclareUnicodeCharacter{1F4A1}{[IDEA]} % Light bulb
% Box-drawing characters (used in ASCII art tables / tree diagrams)
\DeclareUnicodeCharacter{2500}{-} % Horizontal line
\DeclareUnicodeCharacter{2502}{|} % Vertical line
\DeclareUnicodeCharacter{2514}{`} % Up-right corner
\DeclareUnicodeCharacter{251C}{|-} % Tee right
\DeclareUnicodeCharacter{2518}{'} % Up-left corner
\DeclareUnicodeCharacter{250C}{,} % Down-right corner
\DeclareUnicodeCharacter{2510}{.} % Down-left corner
\DeclareUnicodeCharacter{252C}{T} % Tee down
\DeclareUnicodeCharacter{2534}{T} % Tee up
\DeclareUnicodeCharacter{253C}{+} % Cross
% Block / shading
\DeclareUnicodeCharacter{2588}{|} % Full block
\DeclareUnicodeCharacter{2591}{.} % Light shade
\DeclareUnicodeCharacter{2592}{:} % Medium shade
\DeclareUnicodeCharacter{2593}{*} % Dark shade
% Math operators / signs in plain text
\DeclareUnicodeCharacter{2212}{-} % Minus sign (U+2212, NOT hyphen-minus)
\DeclareUnicodeCharacter{2032}{$'$} % Prime
\DeclareUnicodeCharacter{2033}{$''$} % Double prime
\DeclareUnicodeCharacter{2248}{$\approx$} % approximately equal
\DeclareUnicodeCharacter{2245}{$\cong$} % congruent
\DeclareUnicodeCharacter{2026}{\ldots} % horizontal ellipsis
% (U+2192 declared above as rightwards arrow)
% Subscripts and superscripts
\DeclareUnicodeCharacter{2080}{$_0$}
\DeclareUnicodeCharacter{2081}{$_1$}
\DeclareUnicodeCharacter{2082}{$_2$}
\DeclareUnicodeCharacter{2083}{$_3$}
\DeclareUnicodeCharacter{2084}{$_4$}
\DeclareUnicodeCharacter{2085}{$_5$}
\DeclareUnicodeCharacter{2086}{$_6$}
\DeclareUnicodeCharacter{2087}{$_7$}
\DeclareUnicodeCharacter{2088}{$_8$}
\DeclareUnicodeCharacter{2089}{$_9$}
\DeclareUnicodeCharacter{2070}{$^0$}
\DeclareUnicodeCharacter{00B9}{$^1$}
\DeclareUnicodeCharacter{00B2}{$^2$}
\DeclareUnicodeCharacter{00B3}{$^3$}
\DeclareUnicodeCharacter{2074}{$^4$}
\DeclareUnicodeCharacter{2075}{$^5$}
\DeclareUnicodeCharacter{2076}{$^6$}
\DeclareUnicodeCharacter{2077}{$^7$}
\DeclareUnicodeCharacter{2078}{$^8$}
\DeclareUnicodeCharacter{2079}{$^9$}
\DeclareUnicodeCharacter{2717}{\ensuremath{\times}} % Ballot X
\DeclareUnicodeCharacter{2192}{$\to$} % Rightwards arrow
\DeclareUnicodeCharacter{2190}{$\leftarrow$} % Leftwards arrow
\DeclareUnicodeCharacter{2194}{$\leftrightarrow$} % Left-right arrow
\DeclareUnicodeCharacter{2261}{$\equiv$} % identical to
\DeclareUnicodeCharacter{2264}{$\leq$} % less or equal
\DeclareUnicodeCharacter{2265}{$\geq$} % greater or equal
\DeclareUnicodeCharacter{2260}{$\neq$} % not equal
\DeclareUnicodeCharacter{00B1}{$\pm$} % plus-minus
\DeclareUnicodeCharacter{00B5}{$\mu$} % micro / mu
\DeclareUnicodeCharacter{00D7}{$\times$} % multiplication sign
\DeclareUnicodeCharacter{00F7}{$\div$} % division sign
\DeclareUnicodeCharacter{00B0}{$^{\circ}$} % degree sign
\DeclareUnicodeCharacter{2208}{$\in$} % element of
\DeclareUnicodeCharacter{2209}{$\notin$} % not element of
\DeclareUnicodeCharacter{2282}{$\subset$} % subset
\DeclareUnicodeCharacter{2286}{$\subseteq$} % subset or equal
\DeclareUnicodeCharacter{2229}{$\cap$} % intersection
\DeclareUnicodeCharacter{222A}{$\cup$} % union
\DeclareUnicodeCharacter{2205}{$\emptyset$} % empty set
\DeclareUnicodeCharacter{221E}{$\infty$} % infinity
\DeclareUnicodeCharacter{2200}{$\forall$} % for all
\DeclareUnicodeCharacter{2203}{$\exists$} % exists
\DeclareUnicodeCharacter{2202}{$\partial$} % partial derivative
\DeclareUnicodeCharacter{2207}{$\nabla$} % nabla
\DeclareUnicodeCharacter{2211}{$\sum$} % summation
\DeclareUnicodeCharacter{220F}{$\prod$} % product
\DeclareUnicodeCharacter{221A}{$\sqrt{}$} % square root sign
\DeclareUnicodeCharacter{03B1}{$\alpha$} % alpha
\DeclareUnicodeCharacter{03B2}{$\beta$} % beta
\DeclareUnicodeCharacter{03B3}{$\gamma$} % gamma
\DeclareUnicodeCharacter{03B4}{$\delta$} % delta
\DeclareUnicodeCharacter{03B5}{$\varepsilon$} % epsilon
\DeclareUnicodeCharacter{03B6}{$\zeta$} % zeta
\DeclareUnicodeCharacter{03B7}{$\eta$} % eta
\DeclareUnicodeCharacter{03B8}{$\theta$} % theta
\DeclareUnicodeCharacter{03BB}{$\lambda$} % lambda (lowercase)
\DeclareUnicodeCharacter{03C0}{$\pi$} % pi
\DeclareUnicodeCharacter{03C3}{$\sigma$} % sigma
\DeclareUnicodeCharacter{03C4}{$\tau$} % tau
\DeclareUnicodeCharacter{03C6}{$\varphi$} % phi
\DeclareUnicodeCharacter{03C8}{$\psi$} % psi
\DeclareUnicodeCharacter{03C9}{$\omega$} % omega
\DeclareUnicodeCharacter{0394}{$\Delta$} % Delta
\DeclareUnicodeCharacter{03A3}{$\Sigma$} % Sigma
\DeclareUnicodeCharacter{03A8}{$\Psi$} % Psi
\DeclareUnicodeCharacter{03A6}{$\Phi$} % Phi

% ============================================
% COLOR DEFINITIONS
% ============================================

\definecolor{defblue}{RGB}{0,102,204}
\definecolor{thmgreen}{RGB}{34,139,34}
\definecolor{exampleyellow}{RGB}{255,215,0}
\definecolor{warnred}{RGB}{220,20,60}
\definecolor{histpurple}{RGB}{147,112,219}
\definecolor{level1green}{RGB}{144,238,144}
\definecolor{level2yellow}{RGB}{255,255,153}
\definecolor{level3red}{RGB}{255,182,193}

% ============================================
% CUSTOM ENVIRONMENTS - COLOR CODED BOXES
% ============================================
%
% Note (V1.2.1): each box accepts a SINGLE optional argument that is
% interpreted as a FREE-FORM TITLE STRING (appended to the box's
% default title). Earlier revisions tried to forward the optional
% argument as a pgfkey/value list, which caused hundreds of
% "I do not know the key '/tcb/<title>'" errors whenever chapters
% wrote things like \begin{definition}[Blue Boxes]. The current
% definitions wrap the optional argument with title=, so call sites
% that pass either free-form titles OR explicit "title={...}" forms
% both work cleanly.

% Helper: append optional descriptive title after the default box title
\newcommand{\pf@boxtitle}[2]{\ifx\relax#1\relax #2\else #2: #1\fi}

% Blue Box - Definitions
\newtcolorbox{definition}[1][\relax]{
  colback=defblue!5,
  colframe=defblue,
  fonttitle=\bfseries,
  title={\pf@boxtitle{#1}{Definition}},
  breakable,
  enhanced,
  before skip=10pt plus 2pt,
  after skip=10pt plus 2pt,
}

% Green Box - Theorems
\newtcolorbox{theorem}[1][\relax]{
  colback=thmgreen!5,
  colframe=thmgreen,
  fonttitle=\bfseries,
  title={\pf@boxtitle{#1}{Theorem}},
  breakable,
  enhanced,
  before skip=10pt plus 2pt,
  after skip=10pt plus 2pt,
}

% Yellow Box - Examples
\newtcolorbox{example}[1][\relax]{
  colback=exampleyellow!10,
  colframe=exampleyellow!80!black,
  fonttitle=\bfseries,
  title={\pf@boxtitle{#1}{Example}},
  breakable,
  enhanced,
  before skip=10pt plus 2pt,
  after skip=10pt plus 2pt,
}

% Red Box - Warnings
\newtcolorbox{warning}[1][\relax]{
  colback=warnred!5,
  colframe=warnred,
  fonttitle=\bfseries,
  title={\pf@boxtitle{#1}{Warning}},
  breakable,
  enhanced,
  before skip=10pt plus 2pt,
  after skip=10pt plus 2pt,
}

% Purple Box - Historical Context
\newtcolorbox{historical}[1][\relax]{
  colback=histpurple!5,
  colframe=histpurple,
  fonttitle=\bfseries,
  title={\pf@boxtitle{#1}{Historical Context}},
  breakable,
  enhanced,
  before skip=10pt plus 2pt,
  after skip=10pt plus 2pt,
}

\newtcolorbox{historicalnote}[1][\relax]{
  colback=histpurple!5,
  colframe=histpurple,
  fonttitle=\bfseries,
  title={\pf@boxtitle{#1}{Historical Note}},
  breakable,
  enhanced,
}

% (Duplicate Red Box - Warning definition removed: was a verbatim
%  redeclaration of the {warning} tcolorbox already defined at lines
%  142-153 above. pdflatex errors with "Command \warning already
%  defined" otherwise. Removing the duplicate keeps the original
%  definition with its before/after skip parameters intact.)

% Light Green Box - Intuitive Understanding
\newtcolorbox{intuitive}[1][\relax]{
  colback=level1green!15,
  colframe=level1green!80!black,
  fonttitle=\bfseries,
  title={\pf@boxtitle{#1}{Understanding Intuitively}},
  breakable,
  enhanced,
}

% Cyan Box - Key Ideas
\newtcolorbox{keyidea}[1][\relax]{
  colback=cyan!5,
  colframe=cyan!60!black,
  fonttitle=\bfseries,
  title={\pf@boxtitle{#1}{Key Idea}},
  breakable,
  enhanced,
}

% Orange Box - Try It Yourself
\newtcolorbox{tryit}[1][\relax]{
  colback=orange!10,
  colframe=orange!80!black,
  fonttitle=\bfseries,
  title={\pf@boxtitle{#1}{Try It Yourself}},
  breakable,
  enhanced,
}

% Gray Box - Chapter Objectives
\newtcolorbox{chapterobjectives}[1][\relax]{
  colback=gray!10,
  colframe=gray!70!black,
  fonttitle=\bfseries\Large,
  title={\pf@boxtitle{#1}{Chapter Objectives}},
  breakable,
  enhanced,
}

% Level Boxes
\newtcolorbox{level1}[1][\relax]{
  colback=level1green!10,
  colframe=level1green!80!black,
  fonttitle=\bfseries,
  title={\pf@boxtitle{#1}{[L1] Level 1: Intuitive Understanding}},
  breakable,
  enhanced,
  before skip=10pt plus 2pt,
  after skip=10pt plus 2pt,
}

\newtcolorbox{level2}[1][\relax]{
  colback=level2yellow!10,
  colframe=level2yellow!80!black,
  fonttitle=\bfseries,
  title={\pf@boxtitle{#1}{[L2] Level 2: Technical Details}},
  breakable,
  enhanced,
  before skip=10pt plus 2pt,
  after skip=10pt plus 2pt,
}

\newtcolorbox{level3}[1][\relax]{
  colback=level3red!10,
  colframe=level3red!80!black,
  fonttitle=\bfseries,
  title={\pf@boxtitle{#1}{[L3] Level 3: Research \& Verification}},
  breakable,
  enhanced,
  before skip=10pt plus 2pt,
  after skip=10pt plus 2pt,
}

% Advanced box (alias for level3 for backward compatibility)
\newtcolorbox{advanced}[1][\relax]{
  colback=level3red!10,
  colframe=level3red!80!black,
  fonttitle=\bfseries,
  title={\pf@boxtitle{#1}{[L3] Advanced Topic}},
  breakable,
  enhanced,
}

% Backward-compatibility aliases (V1.2.1):
% Earlier chapters used `greenbox` (intended: theorem-style green box)
% and `conjecture` (intended: a tcolorbox styled like a theorem but
% labelled "Conjecture"). Neither was defined in the preamble, causing
% hundreds of "Environment ... undefined" errors. Define them now as
% aliases / siblings of the existing boxes.
\newtcolorbox{greenbox}[1][\relax]{
  colback=thmgreen!5,
  colframe=thmgreen,
  fonttitle=\bfseries,
  title={\pf@boxtitle{#1}{Note}},
  breakable,
  enhanced,
  before skip=10pt plus 2pt,
  after skip=10pt plus 2pt,
}

\newtcolorbox{conjecture}[1][\relax]{
  colback=histpurple!5,
  colframe=histpurple,
  fonttitle=\bfseries,
  title={\pf@boxtitle{#1}{Conjecture}},
  breakable,
  enhanced,
  before skip=10pt plus 2pt,
  after skip=10pt plus 2pt,
}

% Protocol box - experimental / computational protocols
\newtcolorbox{protocol}[1][\relax]{
  colback=cyan!5,
  colframe=cyan!60!black,
  fonttitle=\bfseries,
  title={\pf@boxtitle{#1}{Protocol}},
  breakable,
  enhanced,
  before skip=10pt plus 2pt,
  after skip=10pt plus 2pt,
}

% Problem box - open problems
\newtcolorbox{problem}[1][\relax]{
  colback=orange!5,
  colframe=orange!80!black,
  fonttitle=\bfseries,
  title={\pf@boxtitle{#1}{Open Problem}},
  breakable,
  enhanced,
  before skip=10pt plus 2pt,
  after skip=10pt plus 2pt,
}

% V1.2.1 additions for additional environments referenced across
% several chapters but never defined in the original preamble.
% Each delegates to a tcolorbox styled like the closest semantic kin.
\newtcolorbox{construction}[1][\relax]{
  colback=defblue!5, colframe=defblue, fonttitle=\bfseries,
  title={\pf@boxtitle{#1}{Construction}}, breakable, enhanced,
  before skip=10pt plus 2pt, after skip=10pt plus 2pt,
}
\newtcolorbox{comparison}[1][\relax]{
  colback=gray!10, colframe=gray!70!black, fonttitle=\bfseries,
  title={\pf@boxtitle{#1}{Comparison}}, breakable, enhanced,
  before skip=10pt plus 2pt, after skip=10pt plus 2pt,
}
\newtcolorbox{observation}[1][\relax]{
  colback=cyan!5, colframe=cyan!60!black, fonttitle=\bfseries,
  title={\pf@boxtitle{#1}{Observation}}, breakable, enhanced,
  before skip=10pt plus 2pt, after skip=10pt plus 2pt,
}
\newtcolorbox{chapterabstract}[1][\relax]{
  colback=gray!5, colframe=gray!60!black, fonttitle=\bfseries\Large,
  title={\pf@boxtitle{#1}{Chapter Abstract}}, breakable, enhanced,
  before skip=10pt plus 2pt, after skip=15pt plus 2pt,
}
\newtcolorbox{experiment}[1][\relax]{
  colback=orange!10, colframe=orange!80!black, fonttitle=\bfseries,
  title={\pf@boxtitle{#1}{Experiment}}, breakable, enhanced,
  before skip=10pt plus 2pt, after skip=10pt plus 2pt,
}
\newtcolorbox{principle}[1][\relax]{
  colback=defblue!5, colframe=defblue, fonttitle=\bfseries,
  title={\pf@boxtitle{#1}{Principle}}, breakable, enhanced,
  before skip=10pt plus 2pt, after skip=10pt plus 2pt,
}
\newtcolorbox{mechanism}[1][\relax]{
  colback=cyan!5, colframe=cyan!60!black, fonttitle=\bfseries,
  title={\pf@boxtitle{#1}{Mechanism}}, breakable, enhanced,
  before skip=10pt plus 2pt, after skip=10pt plus 2pt,
}
\newtcolorbox{heuristic}[1][\relax]{
  colback=exampleyellow!10, colframe=exampleyellow!80!black, fonttitle=\bfseries,
  title={\pf@boxtitle{#1}{Heuristic}}, breakable, enhanced,
  before skip=10pt plus 2pt, after skip=10pt plus 2pt,
}
\newtcolorbox{evidence}[1][\relax]{
  colback=thmgreen!5, colframe=thmgreen, fonttitle=\bfseries,
  title={\pf@boxtitle{#1}{Evidence}}, breakable, enhanced,
  before skip=10pt plus 2pt, after skip=10pt plus 2pt,
}
% `prop` is sometimes used in chapters as shorthand for proposition;
% map it to the proposition theorem environment.
\newenvironment{prop}[1][]{\begin{proposition}[#1]}{\end{proposition}}

% ============================================
% CUSTOM COMMANDS
% ============================================

% Margin icons
\newcommand{\iconlevel}[1]{\marginpar{\Large #1}}
\newcommand{\iconintuitive}{\iconlevel{🟢}}
\newcommand{\icontechnical}{\iconlevel{🟡}}
\newcommand{\iconresearch}{\iconlevel{🔴}}
\newcommand{\iconcode}{\iconlevel{💻}}
\newcommand{\icongraph}{\iconlevel{📊}}
\newcommand{\iconimportant}{\iconlevel{🎯}}
\newcommand{\iconwarning}{\iconlevel{⚠️}}

% Mathematical operators and symbols
\newcommand{\R}{\mathbb{R}}
\newcommand{\C}{\mathbb{C}}
\newcommand{\N}{\mathbb{N}}
\newcommand{\Z}{\mathbb{Z}}
\newcommand{\Q}{\mathbb{Q}}

% Principal-branch / multi-valued function operator names
% (declared with \DeclareMathOperator so they typeset upright and
% receive proper spacing in display + inline math).
\DeclareMathOperator{\Log}{Log}     % principal complex logarithm
\DeclareMathOperator{\Arg}{Arg}     % principal argument
\DeclareMathOperator{\Li}{Li}       % polylogarithm
\DeclareMathOperator{\Spec}{Spec}   % spectrum
% physics.sty already defines \Tr, so \DeclareMathOperator errors out and
% aborts the whole book build.  Clear it first, keeping our upright form.
\let\Tr\relax
\DeclareMathOperator{\Tr}{Tr}       % trace
\DeclareMathOperator{\sgn}{sgn}     % sign function
\DeclareMathOperator{\rank}{rank}   % matrix rank
\DeclareMathOperator{\diag}{diag}   % diagonal matrix
% Arithmetic-geometry / BSD / Hodge notation used across Ch 24, 25, 21
\DeclareMathOperator{\ord}{ord}     % order of vanishing
\DeclareMathOperator{\im}{im}       % image
\DeclareMathOperator{\Hdg}{Hdg}     % Hodge classes
\DeclareMathOperator{\Alg}{Alg}     % algebraic classes
\DeclareMathOperator{\Real}{Re}     % real part (long form)
\DeclareMathOperator{\cl}{cl}       % cycle class map
\DeclareMathOperator{\Reg}{Reg}     % regulator
\DeclareMathOperator{\encode}{enc}  % encoding function (P vs NP chapter)
\providecommand{\F}{\mathbb{F}}     % finite field
\providecommand{\CH}{\mathrm{CH}}   % Chow group
\providecommand{\Sha}{\text{\fontencoding{T2A}\selectfont\char88}} % Cyrillic Sha; falls back to "Sh" if font missing
% Safer fallback: define \Sha as a styled "Sh" if Cyrillic not loaded
\renewcommand{\Sha}{\mathrm{Sh}}
\providecommand{\Ch}{\mathrm{Ch}}   % Chern character (when not the symbol)
\providecommand{\ch}{\mathrm{ch}}   % Chern character (lowercase)
\DeclareMathOperator{\Hol}{Hol}     % holonomy
\providecommand{\SU}{\mathrm{SU}}   % special unitary group
\providecommand{\SO}{\mathrm{SO}}   % special orthogonal group
\providecommand{\GL}{\mathrm{GL}}   % general linear group
\providecommand{\Sp}{\mathrm{Sp}}   % symplectic group
\providecommand{\Spin}{\mathrm{Spin}}
\providecommand{\TIME}{\mathrm{TIME}}
\providecommand{\NTIME}{\mathrm{NTIME}}
\providecommand{\SPACE}{\mathrm{SPACE}}
\providecommand{\NSPACE}{\mathrm{NSPACE}}
\providecommand{\BPP}{\mathrm{BPP}}
\providecommand{\BQP}{\mathrm{BQP}}

% Digital sum function
\newcommand{\dthree}[1]{D_3(#1)}

% Fractal Resonance Function
\newcommand{\Rf}[2]{R_f(#1,#2)}

% Consciousness field
\newcommand{\Cfield}{C^{\mu\nu}}

% Resonance coefficient
\newcommand{\rescoeff}[1]{\xi(#1)}

% Omega-space (already defined in LaTeX)
% \Omega is a built-in command

% ============================================
% THEOREM STYLES
% ============================================

\theoremstyle{plain}
\newtheorem{thm}{Theorem}[chapter]
\newtheorem{lemma}[thm]{Lemma}
\newtheorem{proposition}[thm]{Proposition}
\newtheorem{corollary}[thm]{Corollary}

\theoremstyle{definition}
\newtheorem{defn}[thm]{Definition}
\newtheorem{exmp}[thm]{Example}
\newtheorem{exercise}{Exercise}[chapter]

\theoremstyle{remark}
\newtheorem{remark}[thm]{Remark}
\newtheorem{note}[thm]{Note}

% ============================================
% CODE LISTING SETTINGS
% ============================================

\lstset{
  basicstyle=\ttfamily\small,
  keywordstyle=\color{blue}\bfseries,
  commentstyle=\color{gray}\itshape,
  stringstyle=\color{red},
  showstringspaces=false,
  breaklines=true,
  frame=single,
  numbers=left,
  numberstyle=\tiny\color{gray},
  backgroundcolor=\color{gray!10}
}

% Python settings
\lstdefinestyle{python}{
  language=Python,
  morekeywords={numpy,scipy,mpmath}
}

% ============================================
% HYPERREF SETTINGS
% ============================================

\hypersetup{
  colorlinks=true,
  linkcolor=blue,
  citecolor=thmgreen,
  urlcolor=defblue,
  pdftitle={Fractal Resonance Mathematics: The Definitive Textbook},
  pdfauthor={Pablo Cohen},
  pdfsubject={Mathematics, Physics, Consciousness},
  pdfkeywords={Fractal Resonance, Riemann Hypothesis, Consciousness, Quantum Field Theory}
}

% ============================================
% HEADER AND FOOTER
% ============================================

\setlength{\headheight}{14pt} % silence ~45 fancyhdr warnings; required ≥13.6pt
\pagestyle{fancy}
\fancyhf{}
\fancyhead[LE]{\leftmark}
\fancyhead[RO]{\rightmark}
\fancyfoot[C]{\thepage}
\renewcommand{\headrulewidth}{0.4pt}
\renewcommand{\footrulewidth}{0pt}

% ============================================
% TITLE INFORMATION
% ============================================

\title{Fractal Resonance Mathematics\\The Definitive Textbook}
\author{Pablo Cohen\\ORCID: 0009-0002-0734-5565}
\date{\today}

% ============================================
% MISCELLANEOUS
% ============================================

% Line spacing
\onehalfspacing

% Chapter opening customization - Enhanced for v2.0
\usepackage{titlesec}
\titleformat{\chapter}[display]
  {\normalfont\Huge\bfseries\color{pfblue}}
  {\filleft\Huge\color{pfgold}\chaptertitlename\ \thechapter}
  {4ex}
  {\titlerule\vspace{2ex}\filleft}
  [\vspace{2ex}\titlerule]

% Part formatting - Enhanced for v2.0
\titleformat{\part}[display]
  {\normalfont\Huge\bfseries\centering\color{pfblack}}
  {\thepart}
  {20pt}
  {\Huge\color{pfgold}}

% Better page breaking for boxes
\usepackage{needspace}
\newcommand{\boxneedspace}{\needspace{10\baselineskip}}

% Better table spacing
\renewcommand{\arraystretch}{1.3}

% ============================================
% MANUSCRIPT CORRECTION ANNOTATION
% ============================================
%
% \manuscriptcorrection{text} renders a footnote labelled
% "Manuscript Correction (Lean cross-check)" used to flag
% inconsistencies surfaced by Wave 55 Lean refutations.
% The original prose is preserved verbatim above the footnote;
% the correction is documented here so referees can audit
% both the original claim and the corrected value side by side.
\newcommand{\manuscriptcorrection}[1]{%
  \footnote{\textbf{Manuscript Correction (Lean cross-check, Wave 55):} #1}%
}

% ============================================
% EMPIRICAL-CLAIM SCOPE MARKER
% ============================================
%
% \empiricalclaim{citation/dataset detail} renders a footnote
% explicitly marking the preceding numerical claim as an
% empirical / observational result derived from the cited
% dataset, NOT a Lean-formalized theorem. Used to separate
% chapter-level statistical claims (goodness-of-fit, clinical
% accuracy, p-values, etc.) from the project's machine-checked
% Lean content. Added in response to the 2026-06-04 referee
% review (MINOR severity finding: "chapter-level empirical
% claims need explicit scope markers").
\newcommand{\empiricalclaim}[1]{%
  \footnote{\textbf{Empirical/observational claim:} this numerical
  value is derived from the cited dataset, not from a Lean-internal
  derivation. Independent reproduction from primary sources is
  required for referee verification. Dataset / source: #1}%
}

% End of preamble
